% =============================================================================
% Task Description
% =============================================================================
% Source: docs/topic.pdf (thesis topic sheet).
% Bibliography keys: xie2022, adegun2020, li2022survey, xia2020a, bai2021,
% liong2018, zhang2022 (bib/thesis.bib)
% =============================================================================

\thispagestyle{empty}
\chapter*{Task Description}

\noindent\textbf{Improving Micro-Expression Recognition through Temporal and Motion Modeling}

\bigskip

Recognizing facial micro-expressions (MEs) \cite{xie2022} is a particularly difficult
task because these expressions are extremely brief, subtle, and dynamic. Traditional
approaches \cite{adegun2020} based on hand-crafted descriptors often depend on manually
selecting key frames, such as the onset or apex, while deep learning methods
\cite{li2022survey} require large and diverse datasets that are still scarce in this
field. Both directions therefore face limitations, and recognition accuracy remains
modest even on widely used datasets such as CASME-II, SAMM, CAS(ME)$^2$, and SMIC.

Two main factors explain these difficulties. First, accurate temporal localization is
essential, since identifying the exact onset, apex, and offset frames of a
micro-expression is critical for reliable classification. Even small errors in phase
alignment can reduce recognition accuracy \cite{xia2020a}. Second, micro-expressions
involve very weak motion signals, consisting of subtle muscle movements that are easily
confused with neutral expressions, background variations, or camera noise. Prior
studies have attempted to address these challenges through motion magnification to
amplify subtle changes \cite{bai2021}, optical-flow--based motion descriptors
\cite{liong2018}, and spatio-temporal networks \cite{zhang2022}. While these methods
have produced incremental improvements, their impact remains limited, highlighting the
need for a more systematic investigation of temporal and motion modeling strategies.

This thesis will investigate whether improved temporal modeling and motion
representations can enhance recognition of micro-expressions in dynamic facial
expression recognition (FER). The central research question is: how can modeling of
temporal phases and motion cues improve recognition accuracy for subtle and short-lived
expressions? The work will build on established supervised models and evaluate
practical strategies such as temporal attention or hybrid motion representations. The
expected outcome is a contribution that provides a clearer understanding of current
limitations and demonstrates approaches for improving sensitivity to fine-grained
dynamics while maintaining robustness to noise.
